% =====================================================================
% 系统开发工具基础 · 实验报告（第 9--12 题）
%
% 编译方式（在 WSL 中）：
%   latexmk -xelatex -halt-on-error 实验报告.tex
%
% 依赖安装（首次需要）：
%   sudo apt install -y texlive-latex-base texlive-latex-recommended \
%                        texlive-latex-extra latexmk texlive-lang-chinese \
%                        texlive-xetex fonts-noto-cjk
% =====================================================================
\documentclass[UTF8,fontset=ubuntu]{ctexart}

\usepackage{amsmath}
\usepackage{booktabs}
\usepackage{geometry}
\usepackage{listings}
\usepackage[hidelinks]{hyperref}

\geometry{a4paper, margin=2.5cm}

\lstset{
  basicstyle=\ttfamily\small,
  breaklines=true,
  frame=single,
  framerule=0.4pt,
  tabsize=4,
  showstringspaces=false,
  columns=fullflexible,
  keepspaces=true,
}

\title{\textbf{系统开发工具基础 · 实验报告（第 9--12 题）}}
\author{姓名：李云龙 \quad 学号：25030021043}
\date{\today}

\begin{document}

\maketitle

\begin{center}
\begin{tabular}{@{}p{3cm}p{11.5cm}@{}}
\toprule
姓名 & 李云龙 \\
学号 & 25030021043 \\
实验日期 & \today \\
实验环境 & Windows 11 + WSL2（Ubuntu 26.04 LTS）；Python 3.14 + setuptools(build) + pytest 9.1 + PyTorch 2.9.1（CPU） \\
\bottomrule
\end{tabular}
\end{center}

\begin{quote}
\textbf{说明：}本报告中各题“运行结果”均为在 WSL 中实际运行的输出；第 10 题的“失败测试”与第 12 题的损失值等关键过程均如实记录。
\end{quote}

\section{实验目的}
\begin{enumerate}
\item 掌握 Python 包的标准打包流程（pyproject.toml、构建后端、入口点），并在干净虚拟环境中验证 wheel 的可安装性。
\item 掌握“测试驱动 + 编程智能体”的可验证修复循环，以及人工 diff 审查的必要性。
\item 掌握把模糊的协作材料（Issue / 提交信息 / 评审意见）改写成可执行信息的方法。
\item 掌握 PyTorch 线性回归训练循环的完整写法（梯度清零、反向传播、参数更新、评估模式）。
\end{enumerate}

\section{实验九：从源码构建并在干净环境安装 Wheel（q09）}

\subsection{包结构与声明}
\textbf{pyproject.toml}：
\begin{lstlisting}
[build-system]
requires = ["setuptools>=68"]
build-backend = "setuptools.build_meta"

[project]
name = "greetlab-25030021043"
version = "0.1.0"

[project.scripts]
sdt-greet = "greetlab.cli:main"
\end{lstlisting}

\textbf{src/greetlab/cli.py}：
\begin{lstlisting}[language=Python]
import argparse


def main():
    p = argparse.ArgumentParser()
    p.add_argument("--name", required=True)
    a = p.parse_args()
    print(f"Hello, {a.name}!")
\end{lstlisting}

（\lstinline|src/greetlab/__init__.py| 为包标记文件。）

\subsection{构建与安装步骤}
\begin{lstlisting}[language=bash]
cd q09
python3 -m build                                 # 1. 从源码构建
python3 -m venv /tmp/greetenv                    # 2. 新建干净虚拟环境
/tmp/greetenv/bin/pip install dist/*.whl         # 3. 只从 wheel 安装
cd /tmp && /tmp/greetenv/bin/sdt-greet --name 25030021043   # 4. 在 q09 之外运行
\end{lstlisting}

\subsection{运行结果}
\textbf{构建产物}：
\begin{lstlisting}
dist/greetlab_25030021043-0.1.0-py3-none-any.whl
dist/greetlab_25030021043-0.1.0.tar.gz
\end{lstlisting}

\textbf{干净 venv 中只装 wheel 后}：
\begin{lstlisting}
greetlab-25030021043 0.1.0
\end{lstlisting}

\textbf{在 q09 目录之外运行}：
\begin{lstlisting}
$ sdt-greet --name 25030021043
Hello, 25030021043!
\end{lstlisting}

\subsection{要点分析}
\begin{itemize}
\item \lstinline|pyproject.toml| 是标准打包声明：\lstinline|[build-system]| 指定构建后端（setuptools），\lstinline|[project]| 提供元数据，\lstinline|[project.scripts]| 把 \lstinline|greetlab.cli:main| 注册为命令行入口 \lstinline|sdt-greet|；
\item \textbf{干净 venv + 只装 wheel} 模拟真实用户的安装路径，验证 wheel 自包含（不依赖源码目录、不依赖开发环境），这是发布前的“发货检查”；
\item wheel 文件名体现了分发名（连字符归一为下划线）、版本、Python 标签（py3）与平台标签（none-any = 纯 Python 跨平台）。
\end{itemize}

\section{实验十：让编程智能体进入可验证的修复循环（q10）}

\subsection{失败的测试（先写测试）}
\textbf{tests/test\_cli.py}：
\begin{lstlisting}[language=Python]
import sys

import pytest

from greetlab.cli import main


@pytest.mark.parametrize("blank", [" ", "   ", "\t"])
def test_blank_name_exits_with_code_2(monkeypatch, blank):
    monkeypatch.setattr(sys, "argv", ["sdt-greet", "--name", blank])
    with pytest.raises(SystemExit) as excinfo:
        main()
    assert excinfo.value.code == 2
\end{lstlisting}

\subsection{修复前：测试失败}
\begin{lstlisting}
$ python3 -m pytest -q
tests/test_cli.py:11: Failed
----------------------------- Captured stdout call -----------------------------
Hello,  !
FAILED tests/test_cli.py::test_blank_name_exits_with_code_2[ ]
FAILED tests/test_cli.py::test_blank_name_exits_with_code_2[   ]
FAILED tests/test_cli.py::test_blank_name_exits_with_code_2[\t]
3 failed in 0.47s
\end{lstlisting}

\subsection{智能体改动（最小修复）}
\begin{lstlisting}[language=Python]
def main():
    p = argparse.ArgumentParser()
    p.add_argument("--name", required=True)
    a = p.parse_args()
    if not a.name.strip():
        p.error("--name must not be blank")   # argparse 的 error 以 SystemExit(2) 结束
    print(f"Hello, {a.name}!")
\end{lstlisting}

\subsection{修复后：测试通过}
\begin{lstlisting}
$ python3 -m pytest -q
...                                                                      [100%]
3 passed in 0.33s
\end{lstlisting}

\subsection{ai\_log.md（核心提示 / 智能体改动 / 人工验证）}
\begin{lstlisting}
1. 核心提示：sdt-greet --name " "（纯空白）当前仍输出问候语；目标=main 以 SystemExit(2) 结束；
   约束=不改 CLI 语法、不引入依赖；验证命令=python3 -m pytest -q
2. 智能体改动：仅改 src/greetlab/cli.py，print 前加 if not a.name.strip(): p.error(...)
3. 测试演进：修复前 3 个参数化用例失败，修复后 3 passed（红 -> 绿）
4. 人工验证：diff 仅新增 3 行无无关修改；pytest 复跑通过
\end{lstlisting}

\subsection{要点分析}
\begin{itemize}
\item 这是 \textbf{TDD + 智能体闭环}：先用失败测试把“目标”变成可验证的断言（\lstinline|SystemExit(2)|），再让智能体改实现，最后用同一条命令判定成败；
\item \textbf{约束要显式给出}（不改语法、不加依赖），否则智能体可能“顺手”重构出无关改动；
\item 人工 diff 审查是闭环的最后一道闸门：确认改动最小、无副作用后再合并。
\end{itemize}

\section{实验十一：把“无法处理”的协作材料改成可执行信息（q11）}

\subsection{Issue（重写）}
\begin{itemize}
\item \textbf{标题}：\lstinline|sdt-greet| 对空白 \lstinline|--name| 未拒绝，仍输出问候语
\item \textbf{环境}：Linux/WSL，Python 3.14，wheel 安装的 greetlab-25030021043；Windows 上是否复现：待确认
\item \textbf{复现命令}：\lstinline|sdt-greet --name " "|
\item \textbf{期望结果}：拒绝空白姓名，以非零退出码结束
\item \textbf{实际结果}：输出 \lstinline|Hello, !|，退出码 0
\end{itemize}

\subsection{提交信息（重写）}
\textbf{标题}：Reject blank --name in greetlab CLI

\textbf{正文}：传入纯空白字符串时程序仍打印问候语并以 0 退出，会掩盖调用方的参数错误。在 main 中对 \lstinline|a.name| 做 strip 校验，空白时调用 \lstinline|parser.error| 以 SystemExit(2) 拒绝，并补充参数化回归测试。

\subsection{评审意见（重写，标注级别）}
\textbf{[Suggestion]} 行为修复正确，但校验放在 main 内、与解析逻辑分离。建议把校验抽成 \lstinline|--name| 的 \lstinline|type| 函数（如 \lstinline|not_blank|），在参数解析阶段即拒绝，逻辑更内聚、便于直接单测。风险点：现分支已满足要求，但后续新增入口需复用该校验，抽离可避免重复。建议动作：抽取 \lstinline|type| 函数并补参数化测试，下个 PR 完成，不阻塞本次合并。

（全文约 340 字，满足 400 字以内要求。）

\subsection{要点分析}
\begin{itemize}
\item 原始材料“不可执行”：\lstinline|Windows 上运行不了| 缺环境与复现步骤，\lstinline|fix bug| 无信息量，\lstinline|写得不好，重写| 无可操作动作；
\item 重写标准是\textbf{接手人能否直接照做}：Issue 给环境 + 复现 + 期望/实际，未知项显式标注“待确认”；提交信息用祈使句标题 + 问题与方案正文；评审意见给出具体行为、风险、建议动作并标注 Blocking/Suggestion/Nit 级别。
\end{itemize}

\section{实验十二：修复一个可复现的线性回归训练循环（q12）}

\subsection{补全后的代码（q12/train.py）}
\begin{lstlisting}[language=Python]
import torch
from torch import nn

torch.manual_seed(20260907)

x = torch.linspace(-1, 1, 101).reshape(-1, 1)
y = 3 * x - 1

model = nn.Linear(1, 1)
loss_fn = nn.MSELoss()
opt = torch.optim.SGD(model.parameters(), lr=0.1)

for _ in range(200):
    pred = model(x)
    loss = loss_fn(pred, y)
    opt.zero_grad()    # 清空上一轮梯度
    loss.backward()    # 反向传播
    opt.step()         # 更新参数

model.eval()                       # 评估模式
with torch.no_grad():              # 关闭梯度跟踪
    final_loss = loss_fn(model(x), y)
    w = model.weight.item()
    b = model.bias.item()

print(f"final loss = {final_loss.item():.8f}")
print(f"weight = {w:.6f}")
print(f"bias = {b:.6f}")

assert final_loss.item() < 0.001, "final loss 应小于 0.001"
print("OK: final loss < 0.001")
\end{lstlisting}

\subsection{运行结果}
\begin{lstlisting}
$ python3 train.py
final loss = 0.00000000
weight = 2.999997
bias = -1.000000
OK: final loss < 0.001
\end{lstlisting}

\subsection{要点分析}
\begin{itemize}
\item 训练循环三要素缺一不可：\lstinline|zero_grad()| 清空上一轮梯度（否则梯度会累积）、\lstinline|backward()| 反向传播求梯度、\lstinline|step()| 按优化器规则更新参数；
\item 评估阶段用 \lstinline|model.eval()| + \lstinline|torch.no_grad()|：前者切换评估态（影响 Dropout/BatchNorm 等），后者不构建计算图，省内存也更快；
\item 最终 weight $\approx 3$、bias $\approx -1$ 是 200 轮梯度下降\textbf{训练收敛}的结果（\lstinline|torch.manual_seed(20260907)| 保证可复现），并非直接赋值；最终损失 $0 < 0.001$，满足要求。
\end{itemize}

\section{遇到的问题与解决}
\begin{center}
\begin{tabular}{@{}p{5cm}p{5.2cm}p{5cm}@{}}
\toprule
问题 & 原因 & 解决 \\
\midrule
空白姓名未报错 & argparse 只校验参数是否存在，不校验内容 & 在 main 中 strip 后判空并调用 parser.error（退出码 2） \\
测试导入不到 src 布局的包 & src 布局下包不在 sys.path & 增加 conftest.py 把 src 加入 sys.path \\
环境缺少 PyTorch & Ubuntu 26.04 默认未安装 & apt install python3-torch（CPU 版 2.9.1） \\
智能体可能引入无关改动 & 缺少显式约束 & 提示中写明约束（不改语法/不加依赖），并人工 diff 审查 \\
\bottomrule
\end{tabular}
\end{center}

\section{实验总结}
通过本次实验，我完整走通了一个 Python 包从源码到 wheel 再到干净环境安装的发布链路，理解了 pyproject.toml 中构建后端与入口点的作用；实践了“失败测试 -> 智能体修复 -> 人工验证”的可验证修复循环，体会到明确目标、显式约束与 diff 审查三者缺一不可；学会了把模糊的协作材料改写成含环境、复现步骤、期望/实际结果的可执行信息；并补全了 PyTorch 线性回归的完整训练循环，掌握了梯度清零/反传/更新与评估模式的正确用法，最终损失收敛到 0（小于 0.001）。

\section*{附录：文件清单}
\begin{lstlisting}
w3/README.md                  第 9--12 题说明（步骤、原理、实测数据）
w3/q09/pyproject.toml         第 9 题打包声明（包名 greetlab-25030021043）
w3/q09/src/greetlab/          第 9 题包源码（__init__.py、cli.py）
w3/q09/dist/                  第 9 题构建产物（.whl 与 .tar.gz）
w3/q10/src/greetlab/cli.py    第 10 题修复后的实现（空白校验）
w3/q10/tests/test_cli.py      第 10 题失败测试（参数化空白名）
w3/q10/conftest.py            第 10 题测试导入配置
w3/q10/ai_log.md              第 10 题智能体循环记录（<=5 行）
w3/q11/communication.md       第 11 题协作材料改写（<=400 字）
w3/q12/train.py               第 12 题 PyTorch 线性回归训练
实验报告.md / 实验报告.tex / 实验报告.pdf    本报告
\end{lstlisting}

\end{document}
